% RESULTADO_RECUPERADO: reunión local de propietarios ya impresos en el
% capítulo 8. No introduce un axioma ni una versión reducida de APP.
\begin{theorem}[Continuidad conjunta de la holonomía nonádica]
\label{thm:nucleo-continuidad-conexion-nonadica-conjunta}
Sea \(\widehat X_K\) el estado APP--TRIT--TPK enriquecido por hoja,
orientación, residuo, cociente, acarreo, firma, cilindro, frontera, ruta,
memoria y supervivencia. La composición de las nueve relaciones locales
publica una única holonomía
\[
 \Gamma_{9,K}=\mathcal R_{K+8}\circ\cdots\circ\mathcal R_K
 :\widehat X_K\longrightarrow\widehat X_{K+9},
\]
compatible con los truncamientos del sistema inverso. En cada vuelta retorna
la fase y avanza la memoria sin reinicializar el estado; el acarreo conserva
su ley de cociclo, la transferencia actúa simultáneamente sobre sus nueve
componentes y la realización operatorial conserva tanto la componente
ortogonal como el término terminal de la identidad de Stein. Por composición,
estas propiedades determinan secciones compatibles a profundidad arbitraria,
no elecciones independientes nivel a nivel.
\end{theorem}

\begin{proof}
La composición y el retorno sin reinicio son, respectivamente,
\cref{p12-83-c7-s1:def:holonomia-nonadica,p12-83-c7-s1:thm:retorno-nonadico-sin-reinicio}.
La acción conjunta sobre las nueve coordenadas y su transferencia son
\cref{p12-83-c7-s1:thm:funciones-nueve-componentes,p12-83-c7-s1:prop:transferencia-conjunta-nueve}.
La descomposición compresión--expresión y la conservación del terminal son
\cref{p12-83-c7-s1:thm:conexion-compresion-expresion,p12-83-c7-s1:thm:stein-telescopia-terminal}.
La prolongación coinductiva certificada en este mismo capítulo conserva esas
identidades bajo truncamiento; por tanto, la familia compatible pertenece al
límite inverso del estado enriquecido y puede iterarse a cualquier profundidad.
\end{proof}
